\documentclass[../main.tex]{subfiles}
\begin{document}

\section{Theory}
\label{sec:theory}

\subsection{Camouflage deflates the signal}

We first quantify what aggregation does to the class signal.

\begin{proposition}[Signal deflation]
\label{prop:deflation}
Under Definition~\ref{def:ccsbm} with mean aggregation~\eqref{eq:meanagg}, the
expected post-aggregation separation along the discriminant direction is
\begin{equation}
    \delta(\rho) \;=\; \mathbb{E}[t_v \mid y_v = \mathrm{B}]
                     - \mathbb{E}[t_u \mid y_u = \mathrm{H}]
                \;=\; \bigl(1 - \rho - \eta(\rho)\bigr)\,\|\Delta\| ,
    \label{eq:deflation}
\end{equation}
and with the edge-balance identity~\eqref{eq:balance},
$\delta(\rho) = \bigl(1 - (1+c)\rho\bigr)\|\Delta\|$.
\end{proposition}

\begin{proof}
Centre the projection so that $\mathbf{w}^\top\boldsymbol{\mu}_{\mathrm{B}} =
+\|\Delta\|/2$ and $\mathbf{w}^\top\boldsymbol{\mu}_{\mathrm{H}} = -\|\Delta\|/2$.
For a bot $v$, each of its $D$ neighbours is human with probability $\rho$, so by
linearity
$\mathbb{E}[t_v] = (1-\rho)\tfrac{\|\Delta\|}{2} + \rho\bigl(-\tfrac{\|\Delta\|}{2}\bigr)
= \tfrac{\|\Delta\|}{2}(1-2\rho)$.
Symmetrically $\mathbb{E}[t_u] = -\tfrac{\|\Delta\|}{2}(1-2\eta)$ for a human $u$.
Subtracting gives $\tfrac{\|\Delta\|}{2}\bigl[(1-2\rho)+(1-2\eta)\bigr] =
(1-\rho-\eta)\|\Delta\|$. Substituting~\eqref{eq:balance} yields the second form.
\end{proof}

Two consequences are worth stating. The signal vanishes entirely at
$\rho_{\mathrm{null}} = 1/(1+c)$, at which point the two classes have identical
post-aggregation means and \emph{no} classifier acting on aggregated features
alone can do better than chance. And the deflation is linear in $\rho$ with slope
$(1+c)\|\Delta\|$, so higher bot prevalence makes the detector more fragile, not
less---a counterintuitive point, since prevalence is usually treated as a
class-imbalance nuisance rather than a robustness parameter.

The noise term requires more care than the signal term.

\begin{theorem}[Post-aggregation error]
\label{thm:error}
Under Definition~\ref{def:ccsbm}, the projected statistics are asymptotically
Gaussian with separation $\delta(\rho)$ from~\eqref{eq:deflation} and per-class
variances
\begin{equation}
    s_{\mathrm{B}}^2 = \frac{\sigma^2 + \rho(1-\rho)\|\Delta\|^2}{D},
    \qquad
    s_{\mathrm{H}}^2 = \frac{\sigma^2 + \eta(1-\eta)\|\Delta\|^2}{D}.
    \label{eq:variances}
\end{equation}
Defining the post-aggregation signal-to-noise ratio
\begin{equation}
    \mathrm{SNR}(\rho) \;=\; \frac{\delta(\rho)}{s_{\mathrm{B}} + s_{\mathrm{H}}},
    \label{eq:snr}
\end{equation}
the Bayes error of the equal-prior two-class problem is $\Phi(-\mathrm{SNR}(\rho))$,
with equality exact when $s_{\mathrm{B}} = s_{\mathrm{H}}$.
\end{theorem}

\begin{proof}
Condition on a single neighbour of a bot. Its projected feature equals
$\pm\|\Delta\|/2$ according to a Bernoulli$(\rho)$ draw, plus independent
$\mathcal{N}(0,\sigma^2)$ noise. By the law of total variance, the mean component
contributes $\|\Delta\|^2\rho(1-\rho)$ and the noise contributes $\sigma^2$, so each
neighbour has variance $\sigma^2 + \rho(1-\rho)\|\Delta\|^2$. The $D$ neighbours are
independent, so averaging divides the variance by $D$, giving $s_{\mathrm{B}}^2$;
the human case is identical with $\eta$ in place of $\rho$. Asymptotic normality of
$t_v$ follows from the central limit theorem as $D$ grows. For two equal-prior
Gaussians with common standard deviation $s$ and mean gap $\delta$, the
likelihood-ratio test thresholds at the midpoint and attains error
$\Phi(-\delta/2s)$, which is~\eqref{eq:snr} when $s_{\mathrm{B}} = s_{\mathrm{H}} = s$.
\end{proof}

\begin{remark}
The mixture term $\rho(1-\rho)\|\Delta\|^2$ is easy to miss: a calculation that
models only the additive noise obtains $\sigma^2/D$ and therefore
\emph{under}-estimates the post-aggregation error, most severely at intermediate
$\rho$ where the mixture variance peaks. Section~\ref{sec:validation} confirms
the corrected form matches simulation to four decimal places.
\end{remark}

\subsection{The critical camouflage ratio}

We can now answer the deployment question directly.

\begin{corollary}[Critical camouflage ratio]
\label{cor:critical}
Mean aggregation yields a lower Bayes error than the graph-free
classifier~\eqref{eq:rawerr} if and only if $\mathrm{SNR}(\rho) > \|\Delta\|/(2\sigma)$.
Writing $\rho^\star$ for the value of $\rho$ at which equality holds, aggregation
helps precisely when $\rho < \rho^\star$. In the noise-dominated regime
$\sigma^2 \gg \rho(1-\rho)\|\Delta\|^2$ this takes the closed form
\begin{equation}
    \rho^\star \;=\; \frac{1 - D^{-1/2}}{1+c}.
    \label{eq:critical}
\end{equation}
\end{corollary}

\begin{proof}
The ``if and only if'' is immediate from monotonicity of $\Phi$. In the
noise-dominated regime the mixture terms in~\eqref{eq:variances} are negligible, so
$s_{\mathrm{B}} \approx s_{\mathrm{H}} \approx \sigma/\sqrt{D}$ and
$\mathrm{SNR}(\rho) \approx (1-\rho-\eta)\|\Delta\|\sqrt{D}/(2\sigma)$. Setting this
equal to $\|\Delta\|/(2\sigma)$ gives $(1-\rho-\eta)\sqrt{D} = 1$; substituting
$\eta = c\rho$ and solving for $\rho$ yields~\eqref{eq:critical}.
\end{proof}

Equation~\eqref{eq:critical} is the paper's most directly actionable result, and it
says three things. First, at $\rho = 0$ the condition reduces to
$\mathrm{SNR}_{\mathrm{agg}}/\mathrm{SNR}_{\mathrm{raw}} = \sqrt{D}$: aggregation
improves separability by $\sqrt{D}$, recovering the $1/\sqrt{D}$ gain established
by Baranwal et al.~\cite{baranwal2021graph} and confirming our model is calibrated
against the known homophilic baseline. Second, $\rho^\star$ increases with $D$, so
detectors operating on densely connected accounts tolerate more camouflage---the
graph is a genuine asset for hub accounts long after it has become a liability for
sparse ones. Third, $\rho^\star$ decreases in $c$: as bots become more prevalent,
the camouflage budget they need in order to break the detector shrinks.

\begin{remark}[Relation to Ma et al.]
\label{rem:ma}
Equation~\eqref{eq:critical} is not new in the balanced case, and we want to be
explicit about what is and is not ours. For a symmetric CSBM with intra- and
inter-class edge probabilities $p$ and $q$, Ma et al.~\cite[Thm.~2]{ma2022homophily}
prove that a node is better classified after aggregation than before whenever
$\deg(i) > (p+q)^2/(p-q)^2$. Their inter-class neighbour fraction is
$\rho = q/(p+q)$, so $(p-q)/(p+q) = 1-2\rho$ and their condition rearranges to
$\sqrt{D}\,(1-2\rho) > 1$, i.e.\ $\rho < (1 - D^{-1/2})/2$. This is
exactly~\eqref{eq:critical} at $c = 1$. A closed-form threshold for when mean
aggregation is worse than discarding the graph was therefore established at
ICLR~2022, and our corollary is its extension to $c \neq 1$.

What the extension adds is not merely a symbol. Three things change when the classes
are imbalanced. (i) The two classes deflate by \emph{different} amounts, $\rho$ for
bots and $\eta = c\rho$ for humans, so the post-aggregation class midpoint
\emph{moves}; a comparison against a decision boundary held fixed at the raw-feature
midpoint---as in~\cite{ma2022homophily}---is then no longer the right one, which is
precisely the criticism Luan et al.~\cite{luan2023when} raise against that
analysis. Our Theorem~\ref{thm:error} compares Bayes errors instead, and is
insensitive to this. (ii) The edge-balance identity $\eta = c\rho$
in~\eqref{eq:balance} is what couples the two, and it is what makes bot
prevalence a robustness parameter rather than a training-time nuisance. (iii) The
mixture-variance term $\rho(1-\rho)\|\Delta\|^2$ is absent from the covariance
$I/\sqrt{\deg(i)}$ reported in~\cite{ma2022homophily}; including it moves
$\rho^\star$ from $0.53$ to $0.46$ in the configuration of
Remark~\ref{rem:practical}, a $13\%$ correction to the threshold.
\end{remark}

\begin{remark}[Practical reading]
\label{rem:practical}
For a detector with typical $D = 16$ and bot prevalence $\pi_B = 0.3$
($c \approx 0.43$), \eqref{eq:critical} gives $\rho^\star \approx 0.53$ in the
noise-dominated regime and $0.46$ once the mixture variance is included. A bot that
directs roughly half its edges at genuine users has therefore neutralised the graph
entirely. This is not an extreme adversary; it is a modest one, which is why
architectural mitigation is necessary rather than optional.
\end{remark}

\subsection{Spectral edge pruning provably extends the operating regime}
\label{sec:sgtr-theory}

SGTR~\cite{dang2026xhbot} scores each edge by a learned function of the feature
discrepancy $(\mathbf{x}_u - \mathbf{x}_v)^2$, motivated by the per-edge
decomposition of the Dirichlet energy $\mathbf{f}^\top L \mathbf{f} =
\sum_{(u,v)\in\mathcal{E}} (f_u-f_v)^2$. The following lemma shows this quantity is
the right statistic: it separates camouflage from genuine edges by exactly the
squared class gap, independently of the feature dimension.

\begin{lemma}[Energy separability]
\label{lem:energy}
Under Definition~\ref{def:ccsbm}, for a within-class edge $(u,v)$ and a
cross-class (camouflage) edge $(u',v')$,
\begin{equation}
    \mathbb{E}\bigl\|\mathbf{x}_u - \mathbf{x}_v\bigr\|^2 = 2d\sigma^2,
    \qquad
    \mathbb{E}\bigl\|\mathbf{x}_{u'} - \mathbf{x}_{v'}\bigr\|^2 = \|\Delta\|^2 + 2d\sigma^2,
\end{equation}
so the expected energy gap is exactly $\|\Delta\|^2$.
\end{lemma}

\begin{proof}
For a within-class edge both endpoints share a mean, so
$\mathbf{x}_u - \mathbf{x}_v \sim \mathcal{N}(0, 2\sigma^2 I_d)$ and the expected
squared norm is $2d\sigma^2$. For a cross-class edge
$\mathbf{x}_{u'} - \mathbf{x}_{v'} \sim \mathcal{N}(\Delta, 2\sigma^2 I_d)$, and
$\mathbb{E}\|Z\|^2 = \|\mathbb{E}Z\|^2 + \mathrm{tr}\,\mathrm{Cov}(Z)$ gives
$\|\Delta\|^2 + 2d\sigma^2$.
\end{proof}

The gap is a \emph{location} shift that does not grow with $d$, whereas the common
term $2d\sigma^2$ does; the score is therefore informative but increasingly noisy in
high dimension, which is precisely why XHBot learns a projection rather than using
the raw quadratic form. We now quantify the benefit of acting on this score.

\begin{theorem}[SGTR raises the tolerated camouflage]
\label{thm:sgtr}
Suppose soft pruning multiplies camouflage edge weights by $(1-\bar e_c)$ and
within-class edge weights by $(1-\bar e_w)$ in expectation, with
$\bar e_c > \bar e_w$. Then weighted aggregation behaves as unweighted aggregation
at the \emph{effective} camouflage ratio
\begin{equation}
    \rho_{\mathrm{eff}}(\rho)
    \;=\; \frac{\rho(1-\bar e_c)}{\rho(1-\bar e_c) + (1-\rho)(1-\bar e_w)}
    \;<\; \rho .
    \label{eq:rhoeff}
\end{equation}
Consequently, with $\bar e_w = 0$, the detector tolerates true camouflage up to
\begin{equation}
    \rho^\star_{\mathrm{SGTR}}
    \;=\; \frac{\rho^\star}{\rho^\star + (1-\bar e_c)(1-\rho^\star)}
    \;\geq\; \rho^\star,
    \label{eq:rhostarsgtr}
\end{equation}
with equality if and only if $\bar e_c = 0$.
\end{theorem}

\noindent
Equation~\eqref{eq:rhostarsgtr} accounts only for the change in neighbourhood
\emph{composition}. It therefore appears to imply
$\rho^\star_{\mathrm{SGTR}} \to 1$ as $\bar e_c \to 1$: that a sufficiently good
pruner tolerates any camouflage whatsoever. That conclusion is false, and
Section~\ref{sec:achievable} shows why. Down-weighting edges non-uniformly also
reduces the \emph{effective degree}, and hence forfeits part of the $\sqrt{D}$
variance reduction that made aggregation worthwhile in the first place. Once that
cost is included, the ceiling is $1 - 1/D$ rather than $1$, and the operating point
is an interior optimum rather than the corner $\bar e_c \to 1$.

\begin{proof}
Weighted mean aggregation computes $\sum_e w_e \mathbf{x}_{u_e} / \sum_e w_e$. In
expectation a fraction $\rho$ of a bot's edges are camouflage edges carrying weight
$(1-\bar e_c)$ and $(1-\rho)$ are within-class edges carrying $(1-\bar e_w)$;
normalising gives the weight share~\eqref{eq:rhoeff} on human neighbours, which is
exactly the role $\rho$ plays in Proposition~\ref{prop:deflation}. The inequality
$\rho_{\mathrm{eff}} < \rho$ follows since $\bar e_c > \bar e_w$ makes the numerator
shrink faster than the denominator. For~\eqref{eq:rhostarsgtr}, set
$\rho_{\mathrm{eff}}(\rho) = \rho^\star$ with $\bar e_w = 0$ and solve: writing
$\beta = 1-\bar e_c$, the equation $\rho\beta = \rho^\star(\rho\beta + 1 - \rho)$
rearranges to $\rho[\beta(1-\rho^\star) + \rho^\star] = \rho^\star$, giving the
stated form. Monotonicity in $\bar e_c$ and the limits are immediate.
\end{proof}

Theorem~\ref{thm:sgtr} converts an ablation delta into a statement about the
operating regime: a pruning module that attenuates camouflage edges by $75\%$ on
average lifts the tolerated camouflage from $0.53$ to $0.82$ at $D=16$,
$\pi_B = 0.3$. It also clarifies what pruning cannot do. Because
$\rho_{\mathrm{eff}} \to 1$ as $\rho \to 1$ for any fixed $\bar e_c < 1$, no
pruning rule with imperfect discrimination defeats a maximally camouflaged
adversary; SGTR buys operating range, not immunity.

\begin{remark}[What is new here, and what is not]
\label{rem:sgtr-priority}
The \emph{mechanism} in~\eqref{eq:rhoeff}---that edge reweighting acts by rescaling
the effective inter-class edge share, so that weights enter only through the products
of weight and edge rate---is not new. Weiss~\cite[Thm.~2]{weiss2026cost} proves the
corresponding statement for cost-sensitive aggregation with within- and
cross-class weights $w_+ > w_-$, and shows the class-mean direction is preserved iff
$w_+/w_- > q/p$; our~\eqref{eq:rhoeff} is that identity with $w_+ = 1-\bar e_w$ and
$w_- = 1-\bar e_c$. The observation that only a perfect discriminator tolerates
arbitrary heterophily is the contrapositive of the same condition. The related
guarantee for graph attention in~\cite{fountoulakis2023graph} is qualitative but
similar in spirit.

What~\eqref{eq:rhostarsgtr} adds is the \emph{inversion}: composing that rescaling
with the finite-$D$ threshold of Corollary~\ref{cor:critical} to obtain a bound on the
\emph{adversary's} budget as an explicit function of pruner quality $\bar e_c$. This
is a short derivation, and we present it as a corollary of two known ingredients
rather than as an independent discovery. Its value is operational: $\bar e_c$ is
measurable for a deployed pruner, so~\eqref{eq:rhostarsgtr} turns a component
ablation into a quantitative operating-range claim. We note also
that~\cite{weiss2026cost} works in the dense regime where degree effects are absorbed
into $O(1/n)$, which removes the $D^{-1/2}$ term that our threshold depends on; the
finite-$D$ setting is therefore necessary for a statement of this form.
\end{remark}

\subsection{What pruning quality is achievable, and what it costs}
\label{sec:achievable}

Theorem~\ref{thm:sgtr} takes the pruner's quality $\bar e_c$ as given. Two things
are then missing: whether a given $\bar e_c$ is attainable at all by a
feature-based score, and what is given up in exchange. We supply both, which
removes the free parameter and corrects the limit noted above.

\paragraph{Only the retained-weight ratio matters.}
Write $a = 1-\bar e_c$ and $b = 1-\bar e_w$ for the surviving weight on camouflage
and within-class edges. Weighted mean aggregation is invariant to a common rescaling
of all weights, so $a$ and $b$ enter only through the ratio
\begin{equation}
    \theta \;=\; a/b \;\in\; (0,1] ,
    \label{eq:theta}
\end{equation}
which we call the \emph{retained-weight ratio}; $\theta = 1$ is no pruning and
$\theta = 0$ is perfect pruning. This is the same quantity that governs the
cost-sensitive condition $w_+/w_- > q/p$ of~\cite{weiss2026cost}. In terms of
$\theta$, the effective camouflage ratio of Theorem~\ref{thm:sgtr} is
$\rho_{\mathrm{eff}} = \rho\theta/(\rho\theta + 1-\rho)$.

\begin{proposition}[Pruning forfeits effective degree]
\label{prop:deff}
A weighted mean of $D$ unit-variance terms has the variance of an unweighted mean
of
\begin{equation}
    D_{\mathrm{eff}}
    \;=\; D\,\frac{\bigl[(1-\rho) + \rho\theta\bigr]^2}
                   {(1-\rho) + \rho\theta^{2}}
    \;\leq\; D
    \label{eq:deff}
\end{equation}
terms, with equality if and only if $\theta = 1$. Consequently pruned aggregation
beats the graph-free classifier if and only if
\begin{equation}
    \bigl|1 - (1+c)\,\rho_{\mathrm{eff}}(\rho,\theta)\bigr|
    \sqrt{D_{\mathrm{eff}}(\rho,\theta)} \;>\; 1 .
    \label{eq:pruned-condition}
\end{equation}
\end{proposition}

\begin{proof}
The variance of $\sum_e w_e z_e / \sum_e w_e$ for independent unit-variance $z_e$ is
$\sum_e w_e^2/(\sum_e w_e)^2$, so the variance-equivalent count is
$(\sum_e w_e)^2/\sum_e w_e^2$. Substituting $\rho D$ edges of weight $a$ and
$(1-\rho)D$ of weight $b$, and cancelling $b^2$, gives~\eqref{eq:deff}; the
inequality is Cauchy--Schwarz, with equality iff the weights are constant.
Combining with Proposition~\ref{prop:deflation} applied at $\rho_{\mathrm{eff}}$
yields~\eqref{eq:pruned-condition}.
\end{proof}

This resolves a degeneracy that~\eqref{eq:rhostarsgtr} alone does not. Uniform
attenuation ($\theta = 1$, any common scale) leaves both $\rho_{\mathrm{eff}}$ and
$D_{\mathrm{eff}}$ unchanged and so achieves exactly nothing, as it must; and
aggressive pruning is now penalised rather than free.

\begin{corollary}[Ceiling for a perfect pruner]
\label{cor:ceiling}
Even with $\theta = 0$---every camouflage edge removed, every genuine edge
retained---aggregation beats the graph-free classifier if and only if
\begin{equation}
    \rho \;<\; 1 - \frac{1}{D} .
    \label{eq:ceiling}
\end{equation}
\end{corollary}

\begin{proof}
At $\theta = 0$ we have $\rho_{\mathrm{eff}} = 0$, so the bias factor
in~\eqref{eq:pruned-condition} is $1$, and
$D_{\mathrm{eff}} = D(1-\rho)^2/(1-\rho) = D(1-\rho)$. The condition becomes
$\sqrt{D(1-\rho)} > 1$, i.e.\ $\rho < 1 - 1/D$.
\end{proof}

The interpretation is worth stating plainly, because it is the opposite of what
Theorem~\ref{thm:sgtr} suggests in isolation. A perfect pruner does not confer
immunity. It removes the bias but leaves a bot with only $(1-\rho)D$ usable
neighbours, and once that number falls to one the graph carries no more information
than the node's own features. Pruning cannot manufacture evidence that camouflage
has destroyed; at $\rho \to 1$ there is simply nothing left to aggregate.

\paragraph{Achievable $\theta$.}
It remains to bound $\theta$ from below. SGTR scores an edge by the feature
discrepancy $s_{uv} = \|\mathbf{x}_u - \mathbf{x}_v\|^2$, whose law under the CSBM
is exactly
\begin{equation}
    s \sim 2\sigma^2\chi^2_{d}
    \quad\text{(within class)},
    \qquad
    s \sim 2\sigma^2\chi^2_{d}\!\left(\tfrac{r^2}{2}\right)
    \quad\text{(camouflage)},
    \label{eq:scorelaw}
\end{equation}
with $r = \|\Delta\|/\sigma$ and $\chi^2_d(\lambda)$ the noncentral chi-square. Any
threshold rule therefore has an operating point $(\bar e_w, \bar e_c)$ on the ROC
of~\eqref{eq:scorelaw}, and $\theta$ is the ratio of its complements. In the
large-$d$ Gaussian approximation the discriminability index is
\begin{equation}
    \kappa \;=\; \frac{\|\Delta\|^2}{2\sigma^2\sqrt{2d}}
            \;=\; \frac{r^2}{2\sqrt{2d}} ,
    \label{eq:kappa}
\end{equation}
so $\bar e_c \approx \Phi(\Phi^{-1}(\bar e_w) + \kappa)$.

Two consequences follow. First, $\bar e_w \to 0$ forces $\bar e_c \to 0$: there is no
threshold that removes camouflage edges without also removing genuine ones, so the
idealisation $\bar e_w = 0$ used in~\eqref{eq:rhostarsgtr} is not attainable and the
joint form is required. Second, because $\kappa$ scales as $d^{-1/2}$, the raw
discrepancy score \emph{degrades} as the feature dimension grows---the two score
distributions concentrate and overlap. This is the formal content of the remark that
motivates learning a projection: mapping to $k \ll d$ dimensions while preserving
$\|\Delta\|$ multiplies $\kappa$ by $\sqrt{d/k}$.

Substituting the achievable $\theta$ into~\eqref{eq:pruned-condition} gives a
tolerance that depends only on $(D, c, r, d)$---all measurable---with no free pruning
parameter. Section~\ref{sec:validation} reports the resulting values.

\subsection{Depth}
\label{sec:depth}

Everything so far concerns a single layer. Since deployed detectors stack two or
more, we record how the picture changes. Assume the graph is locally tree-like, so
that the neighbourhoods aggregated at different levels are disjoint; this is the
standard local-weak-convergence approximation and is accurate for sparse graphs in
which short cycles are rare.

\begin{theorem}[Depth recursion]
\label{thm:depth}
Let $\Delta_l$ denote the class-mean separation and $V_B^{(l)}, V_H^{(l)}$ the
class-conditional variances after $l$ layers, with
$\Delta_0 = \|\Delta\|$ and $V^{(0)}_B = V^{(0)}_H = \sigma^2$. Then
\begin{align}
    \Delta_{l+1} &= \lambda\,\Delta_l, \qquad \lambda = 1-(1+c)\rho,
    \label{eq:depth-mean}\\
    V^{(l+1)}_B &= \tfrac{1}{D}\Bigl[(1-\rho)V^{(l)}_B + \rho V^{(l)}_H
                    + \rho(1-\rho)\Delta_l^2\Bigr],
    \label{eq:depth-varB}\\
    V^{(l+1)}_H &= \tfrac{1}{D}\Bigl[\eta V^{(l)}_B + (1-\eta)V^{(l)}_H
                    + \eta(1-\eta)\Delta_l^2\Bigr].
    \label{eq:depth-varH}
\end{align}
In the noise-dominated regime, where the $\Delta_l^2$ terms are negligible, this
collapses to
\begin{equation}
    \frac{\mathrm{SNR}_L}{\mathrm{SNR}_{\mathrm{raw}}}
    \;=\; \bigl(\lambda\sqrt{D}\bigr)^{L}.
    \label{eq:depth-clean}
\end{equation}
\end{theorem}

\begin{proof}
Deferred to Appendix~\ref{app:depth}.
\end{proof}

Three consequences. First, \textbf{the crossover is depth-independent}: since
$(\lambda\sqrt D)^L > 1$ iff $\lambda\sqrt D > 1$ for every $L \geq 1$, the critical
ratio is the same at every depth, and $\rho^\star$ is not an artefact of the
single-layer analysis. Second, \textbf{depth amplifies the consequence
exponentially}: below $\rho^\star$ each additional layer multiplies the SNR by
$\lambda\sqrt D > 1$, and above it by $\lambda\sqrt D < 1$. Stacking layers is
therefore a leveraged bet on being on the right side of the threshold, which is a
sharper statement than the usual observation that deep GNNs over-smooth: the sign of
the effect is set by the topology, not by the depth.

Third, the harmful region is \emph{bounded above}. Since only $|\lambda|$ matters,
aggregation is harmful precisely on the window
\begin{equation}
    \rho \;\in\;
    \left(\frac{1 - D^{-1/2}}{1+c},\;\; \frac{1 + D^{-1/2}}{1+c}\right),
    \label{eq:window}
\end{equation}
and \emph{helps again} above it. The reason is instructive: a bot with
$\rho \to 1$ connects almost exclusively to humans, so its neighbourhood is strongly
\emph{anti}-correlated with its label and therefore informative once more, with the
orientation reversed---a downstream linear layer can absorb the sign. This is the
adversarial instance of the observation that heterophily is not inherently
harmful~\cite{ma2022homophily,luan2023when}, and it identifies the genuinely
dangerous strategy as \emph{intermediate} camouflage, which maximises ambiguity
rather than maximising heterophily. For $D=16$ and $\pi_B = 0.3$ the window is
$(0.525, 0.875)$.

\subsection{Why a non-mean channel is needed}

The expressivity argument for tri-channel aggregation is a direct application of the
standard multiset-injectivity analysis of Xu et al.~\cite{xu2019powerful}, and we
record it as a proposition rather than claim it as a new result. It is included
because it settles the third question of Section~\ref{sec:intro}---whether
multi-channel aggregation is more \emph{capable} or merely better
parameterised---which an ablation cannot settle.

\begin{proposition}[Expressivity separation; after~\cite{xu2019powerful}]
\label{thm:thca}
Let $\mathcal{F}_{\mathrm{mean}}$ be the class of functions computable by any
single-channel homophilic aggregator of the form
$\mathbf{h}_v = \phi\bigl(W \cdot \mathrm{mean}_{u\in\mathcal{N}(v)} \mathbf{x}_u\bigr)$,
and $\mathcal{F}_{\mathrm{THCA}}$ the class computable by tri-channel aggregation
with homophilic, max-pooling, and self channels under attention fusion. Then
$\mathcal{F}_{\mathrm{mean}} \subsetneq \mathcal{F}_{\mathrm{THCA}}$.
\end{proposition}

\begin{proof}
\emph{Containment.} Setting the fusion weights to $\alpha_H = 1$,
$\alpha_X = \alpha_S = 0$ reproduces the homophilic channel exactly, so every
$f \in \mathcal{F}_{\mathrm{mean}}$ is realised.

\emph{Strictness.} Consider the scalar neighbourhood multisets
$\mathcal{N}_1 = \{0, 2\}$ and $\mathcal{N}_2 = \{1, 1\}$. Both have mean $1$, so
$\mathrm{mean}(\mathcal{N}_1) = \mathrm{mean}(\mathcal{N}_2)$ and hence
$\phi(W\cdot 1)$ is identical for both under \emph{every} choice of $\phi$ and $W$:
no member of $\mathcal{F}_{\mathrm{mean}}$ separates them. The max-pooling channel
with identity $\mathrm{MLP}_X$ returns $\max(\mathcal{N}_1) = 2 \neq 1 =
\max(\mathcal{N}_2)$, so a tri-channel aggregator with $\alpha_X > 0$ assigns them
distinct representations. Hence the inclusion is strict.
\end{proof}

The separating instance is not a pathology: $\{0,2\}$ versus $\{1,1\}$ is precisely
the contrast between a node with one extreme neighbour and one with two typical
neighbours---the signature of an account with a small number of highly anomalous
connections. Mean aggregation is blind to it by construction, which is the
structural reason a heterophilic channel is needed rather than merely helpful.

\subsection{What the disentanglement objective does and does not guarantee}
\label{sec:cpd-theory}

CPD~\cite{dang2026xhbot} combines a supervised prototype InfoNCE term with an
orthogonality penalty $\mathcal{L}_{MI} = \frac{1}{N}\sum_v \|(\mathbf{z}^b_v)^\top
\mathbf{z}^s_v\|_F^2$. The two terms warrant different levels of confidence, and we
state them separately.

\begin{proposition}[The prototype term maximises an MI lower bound]
\label{prop:nce}
Let $\mathcal{L}_{NCE}$ be the supervised prototype InfoNCE loss over $K$ class
prototypes. Then
\begin{equation}
    I(\mathbf{z}^b; y) \;\geq\; \log K - \mathcal{L}_{NCE},
\end{equation}
so minimising $\mathcal{L}_{NCE}$ maximises a variational lower bound on the mutual
information between the behavioural code and the label.
\end{proposition}

\begin{proof}
This is the InfoNCE bound of~\cite{oord2018representation} in the form given by
Poole et al.~\cite{poole2019variational}, instantiated with the $K$ prototypes as
the contrastive set and the critic $f(\mathbf{z}, k) = (\mathbf{z}^b)^\top
\mathbf{P}_k / \tau_{nce}$.
\end{proof}

\begin{remark}
The bound saturates at $\log K$, which for binary bot detection is $\log 2$. The
objective therefore cannot certify more than one bit of label information, and its
practical value lies in shaping the geometry of the code rather than in the
tightness of the bound---a caveat we consider important to state, since the
$\log K$ ceiling is often left implicit.
\end{remark}

\begin{proposition}[The orthogonality penalty controls dependence only under a
Gaussian assumption]
\label{prop:mi}
If $(\mathbf{z}^b, \mathbf{z}^s)$ are jointly Gaussian with cross-correlation $r$,
then $I(\mathbf{z}^b;\mathbf{z}^s) = -\tfrac{1}{2}\log(1-r^2)$, so
$\mathcal{L}_{MI} \to 0$ implies $I \to 0$. Without the joint-Gaussian assumption
the implication fails: $\mathcal{L}_{MI} = 0$ enforces only uncorrelatedness, and
there exist uncorrelated but strongly dependent codes.
\end{proposition}

\begin{proof}
The identity for the bivariate Gaussian is standard. For the converse, take
$z^b \sim \mathcal{N}(0,1)$ and $z^s = (z^b)^2 - 1$: then
$\mathbb{E}[z^b z^s] = \mathbb{E}[(z^b)^3] - \mathbb{E}[z^b] = 0$, so the penalty
vanishes, yet $z^s$ is a deterministic function of $z^b$ and
$I(z^b; z^s) = \infty$ in the continuous sense.
\end{proof}

We draw the honest conclusion. $\mathcal{L}_{MI}$ is a \emph{surrogate}, exact only
in the Gaussian case, and the empirical success of CPD should not be read as
evidence of guaranteed independence between behavioural and positional codes.
Moreover, by the impossibility result of Locatello
et al.~\cite{locatello2019challenging}, no purely unsupervised objective can
identify the two factors; what makes the behavioural code identifiable here is the
\emph{label supervision} in the prototype term, which anchors one of the two
subspaces. The positional code carries no such anchor and is identified only up to
the residual subspace---a limitation of the design, not of the analysis.

\subsection{Computational cost}

\begin{proposition}[No asymptotic overhead]
\label{prop:complexity}
For a graph with $N$ nodes, $|\mathcal{E}|$ edges, and hidden width $d$, one layer
of SGTR followed by tri-channel aggregation and attention fusion costs
$O(|\mathcal{E}|d + Nd^2)$ time and $O(|\mathcal{E}| + Nd)$ memory---the same order
as a single relational graph convolution.
\end{proposition}

\begin{proof}
SGTR evaluates a fixed-width scoring function per edge: $O(|\mathcal{E}|d)$. The
homophilic and max-pooling channels each perform one pass over the edge set with a
$d \times d$ transform per node, $O(|\mathcal{E}|d + Nd^2)$; the self channel is
$O(Nd^2)$; attention fusion computes three scalars per node, $O(Nd)$. Summing the
three channels multiplies the constant by a factor of at most three and leaves the
order unchanged. Memory is dominated by storing edge weights and node states.
\end{proof}

The practical content is that the theoretical robustness gains of
Theorem~\ref{thm:sgtr} and Theorem~\ref{thm:thca} are obtained without changing the
asymptotic cost profile, which is consistent with the sub-millisecond per-node
inference reported for the deployed system~\cite{dang2026xhbot}.

\end{document}
